% Gliederung Stand 30.10.2026. ~14 Seiten. Nur übertragbare Details, keine Entwicklungsgeschichte, keine Bedienschritte.
\chapter{Umsetzung}
\label{chap:umsetzung}
% Kurze Kapiteleinleitung vor 7.1.

\section{Containerisierte Bereitstellung}
\label{sec:bereitstellung}
% Images für Sim und Real, Profile, Startmenü als Einstieg. [A6]

\section{Simulationsumgebung}
\label{sec:simulationsumgebung}
% [A3, A4]

\subsection{Kopplung von MuJoCo und DDS}
\label{sec:bridge}
% Bridge, Befehls-Merge lowcmd + arm_sdk.

\subsection{Zeitschritt und Echtzeitbetrieb}
\label{sec:zeitschritt}
% 1 kHz, Lockstep, Echtzeitfaktor.

\subsection{Szenen und Arbeitsplatzmodelle}
\label{sec:szenen}
% Szenen-Editor, Greifobjekte, Störungen.

\section{Fortbewegung}
\label{sec:fortbewegung_umsetzung}

\subsection{Ersatzregler in der Simulation}
\label{sec:ersatzregler}
% Stand-PD mit Schwerpunktführung, Lauf-Policies, Übergabelogik.

\subsection{Anbindung des Onboard-Reglers}
\label{sec:onboard_anbindung}

\section{Armregelung und Bahnplanung}
\label{sec:armregelung}

\subsection{Regelkette und inverse Kinematik}
\label{sec:regelkette}
% 250-Hz-Takt, IK in eigenem Prozess.

\subsection{Kollisionsmodell und Laufzeitüberwachung}
\label{sec:kollisionsmodell}

\subsection{Bahnplaner}
\label{sec:bahnplaner}
% Trichter-Abstandsmodell, Rückfallstufen.

\section{Inspire-Hände}
\label{sec:haende}
% Einheitliche Schnittstelle mit drei Backends, Kraftregelung, Taktilzonen.

\section{Bedienwege}
\label{sec:bedienwege}
% [F3] Jeweils Konzept + Begründung, keine Bedienschritte.

Tabelle~\ref{tab:bedienwege} \pruefen{Entwurfstabelle aus dem früheren Ergebnis-Kapitel, Inhalt und Platzierung prüfen}.

\begin{table}[htbp]
    \centering
    \caption{Verschiedene Arten der Bedienung}
    \label{tab:bedienwege}
    \begin{tabular}{|l|l|l|l|l|l|}\hline
           &Demo-GUI& Rviz&Controller&  Arm-API &Handy-GUI\\\hline
           Arme&x& x&x& x&\\\hline
     Beine& x& Auto-Nav& x& &\\\hline
     Hände& x& & x& x&x\\\hline
     Posenspeicher& x& & & x&x\\ \hline
    \end{tabular}
\end{table}

\subsection{Demo-Oberfläche}
\label{sec:demo_gui}

\subsection{Gamepad}
\label{sec:gamepad}
% PS4/Xbox.

\subsection{Smartphone}
\label{sec:smartphone}
% Bluetooth.

\subsection{Entwickler-Oberfläche und Arm-API}
\label{sec:entwickler_api}
% Streamdeck, HTTP/JSON-Schnittstelle.

\section{Navigation in der Simulation}
\label{sec:navigation}
% Kurz: aus g1pilot übernommen und eingebunden, Beispiel für ein integriertes Modul.
